\documentclass[11pt]{article}
\newcommand{\doctitle}{Implementing Direct Mean Moderation for
  Biomarker-Treatment Interactions:
  A Worked Example in Signal Calibration}
\newcommand{\shorttitle}{Mean Moderation Implementation}
\newcommand{\docdate}{2026-03-25}
% pmsimstats-ng standard document preamble
% Usage: % pmsimstats-ng standard document preamble
% Usage: % pmsimstats-ng standard document preamble
% Usage: \input{pmsimstats-preamble} after \documentclass[11pt]{article}
%
% Documents must define before \input:
%   \newcommand{\doctitle}{...}       Full document title
%   \newcommand{\shorttitle}{...}     Short title for header
%   \newcommand{\docdate}{YYYY-MM-DD} Document date

\usepackage[margin=1in]{geometry}
\usepackage{amsmath,amssymb}
\usepackage[T1]{fontenc}
\usepackage{lmodern}
\usepackage{microtype}
\usepackage{graphicx}
\graphicspath{{figures/}}
\usepackage{booktabs}
\usepackage{longtable}
\usepackage{array}
\usepackage{hyperref}
\usepackage{xcolor}
\usepackage{fancyhdr}
\usepackage{parskip}
\usepackage{caption}
\usepackage{enumitem}
\usepackage{verbatim}
\usepackage{listings}

\lstset{
  language=R,
  basicstyle=\small\ttfamily,
  keywordstyle=\color{blue!70!black},
  commentstyle=\color{green!50!black},
  stringstyle=\color{red!60!black},
  breaklines=true,
  frame=single,
  framerule=0.5pt,
  rulecolor=\color{gray!40},
  backgroundcolor=\color{gray!5},
  xleftmargin=1em,
  framexleftmargin=1em,
  columns=flexible
}

\hypersetup{
  colorlinks=true,
  linkcolor=blue!60!black,
  citecolor=green!50!black,
  urlcolor=blue!60!black
}

\pagestyle{fancy}
\fancyhf{}
\lhead{\footnotesize pmsimstats-ng Technical Report}
\rhead{\footnotesize \shorttitle}
\rfoot{\footnotesize \thepage}
\renewcommand{\headrulewidth}{0.4pt}

\title{\doctitle}
\author{pmsimstats team}
\date{\docdate}
 after \documentclass[11pt]{article}
%
% Documents must define before \input:
%   \newcommand{\doctitle}{...}       Full document title
%   \newcommand{\shorttitle}{...}     Short title for header
%   \newcommand{\docdate}{YYYY-MM-DD} Document date

\usepackage[margin=1in]{geometry}
\usepackage{amsmath,amssymb}
\usepackage[T1]{fontenc}
\usepackage{lmodern}
\usepackage{microtype}
\usepackage{graphicx}
\graphicspath{{figures/}}
\usepackage{booktabs}
\usepackage{longtable}
\usepackage{array}
\usepackage{hyperref}
\usepackage{xcolor}
\usepackage{fancyhdr}
\usepackage{parskip}
\usepackage{caption}
\usepackage{enumitem}
\usepackage{verbatim}
\usepackage{listings}

\lstset{
  language=R,
  basicstyle=\small\ttfamily,
  keywordstyle=\color{blue!70!black},
  commentstyle=\color{green!50!black},
  stringstyle=\color{red!60!black},
  breaklines=true,
  frame=single,
  framerule=0.5pt,
  rulecolor=\color{gray!40},
  backgroundcolor=\color{gray!5},
  xleftmargin=1em,
  framexleftmargin=1em,
  columns=flexible
}

\hypersetup{
  colorlinks=true,
  linkcolor=blue!60!black,
  citecolor=green!50!black,
  urlcolor=blue!60!black
}

\pagestyle{fancy}
\fancyhf{}
\lhead{\footnotesize pmsimstats-ng Technical Report}
\rhead{\footnotesize \shorttitle}
\rfoot{\footnotesize \thepage}
\renewcommand{\headrulewidth}{0.4pt}

\title{\doctitle}
\author{pmsimstats team}
\date{\docdate}
 after \documentclass[11pt]{article}
%
% Documents must define before \input:
%   \newcommand{\doctitle}{...}       Full document title
%   \newcommand{\shorttitle}{...}     Short title for header
%   \newcommand{\docdate}{YYYY-MM-DD} Document date

\usepackage[margin=1in]{geometry}
\usepackage{amsmath,amssymb}
\usepackage[T1]{fontenc}
\usepackage{lmodern}
\usepackage{microtype}
\usepackage{graphicx}
\graphicspath{{figures/}}
\usepackage{booktabs}
\usepackage{longtable}
\usepackage{array}
\usepackage{hyperref}
\usepackage{xcolor}
\usepackage{fancyhdr}
\usepackage{parskip}
\usepackage{caption}
\usepackage{enumitem}
\usepackage{verbatim}
\usepackage{listings}

\lstset{
  language=R,
  basicstyle=\small\ttfamily,
  keywordstyle=\color{blue!70!black},
  commentstyle=\color{green!50!black},
  stringstyle=\color{red!60!black},
  breaklines=true,
  frame=single,
  framerule=0.5pt,
  rulecolor=\color{gray!40},
  backgroundcolor=\color{gray!5},
  xleftmargin=1em,
  framexleftmargin=1em,
  columns=flexible
}

\hypersetup{
  colorlinks=true,
  linkcolor=blue!60!black,
  citecolor=green!50!black,
  urlcolor=blue!60!black
}

\pagestyle{fancy}
\fancyhf{}
\lhead{\footnotesize pmsimstats-ng Technical Report}
\rhead{\footnotesize \shorttitle}
\rfoot{\footnotesize \thepage}
\renewcommand{\headrulewidth}{0.4pt}

\title{\doctitle}
\author{pmsimstats team}
\date{\docdate}


\begin{document}
\maketitle
\thispagestyle{fancy}

\begin{abstract}
This document traces four attempts to implement direct mean
moderation (Architecture~A) for the biomarker-treatment
interaction in the \texttt{pmsimstats-ng} clinical trial
simulation framework. Three attempts produced incorrect
behavior due to scaling errors: one generated no signal during
off-drug periods, another generated signal that decayed too
rapidly, and a third generated signal so large that power
reached 1.00 under all conditions. The correct formulation
uses standardized biomarker values, a binary on-drug
indicator, and scaling by the biological response standard
deviation to produce effect sizes that are calibrated against
the existing multivariate normal (MVN) architecture. We
present the mathematical derivation, explain why each
incorrect attempt failed, and provide empirical verification.
The narrative is intended to illustrate how dimensional
analysis and calibration against a known reference
implementation can guide the development of simulation
components.
\end{abstract}

\tableofcontents

\section{Introduction}

\subsection{The problem}

Consider a clinical trial in which a baseline biomarker
$B_i$ (e.g., resting blood pressure) is hypothesized to
predict which participants will benefit from a drug. A
simulation framework generates synthetic trial data and
estimates the statistical power to detect this
biomarker-treatment interaction. How should the simulation
encode the relationship between biomarker and drug response?

There are two fundamentally different approaches, which we
call Architecture~A and Architecture~B. This document is
concerned with the implementation of Architecture~A.

\subsection{Architecture B: differential correlation (MVN)}

The existing implementation in \texttt{pmsimstats-ng}
generates all participant-level variables---biomarker ($B_i$),
baseline symptom score ($\text{BL}_i$), and response
components at each timepoint ($\text{BR}_{it}$,
$\text{TV}_{it}$, $\text{PB}_{it}$)---from a single
multivariate normal (MVN) draw:
%
\begin{equation}
  \mathbf{X}_i \sim \text{MVN}(\boldsymbol{\mu},\,\Sigma)
\end{equation}
%
The biomarker-treatment interaction is encoded in the
covariance matrix $\Sigma$. Specifically, the correlation
between the biomarker and the biological response (BR)
component depends on treatment status:
%
\begin{equation}\label{eq:mvn-cor}
  \text{Cor}(B_i,\,\text{BR}_{it}) =
  \begin{cases}
    c_{bm} & \text{if on drug at time } t \\
    c_{bm} \cdot e^{-\lambda \cdot t_{sd}} &
      \text{if off drug (with decay)} \\
    0 & \text{if never treated}
  \end{cases}
\end{equation}
%
where $c_{bm}$ is the biomarker-BR correlation parameter,
$\lambda$ is the correlation decay rate, and $t_{sd}$ is time
since drug discontinuation.

The interaction signal lives in the \emph{second moment}
(covariance) of the joint distribution. It is not visible in
the population means---all participants have the same expected
BR at a given timepoint regardless of their biomarker value.
The signal emerges only in the conditional distribution:
given that participant $i$ has biomarker value $b$, their
expected BR is shifted:
%
\begin{equation}\label{eq:cond-exp}
  E[\text{BR}_{it} \mid B_i = b,\,\text{on drug}]
  = \mu_{\text{BR}}(t) + c_{bm} \cdot
    \frac{\sigma_{\text{BR}}}{\sigma_B}
    \cdot (b - \mu_B)
\end{equation}
%
This is a standard property of the multivariate normal
distribution (see, e.g., Mardia et al., 1979, Chapter~3).

\subsection{Architecture A: direct mean moderation}

Under Architecture~A, the biomarker enters the
\emph{mean structure} directly. The BR component for
participant $i$ at timepoint $t$ is:
%
\begin{equation}\label{eq:arch-a}
  \text{BR}_{it} = \text{BR}_{it}^{\text{base}}
  + \beta_{bm} \cdot B_i^{*} \cdot D_{it} \cdot \sigma_{\text{BR}}
\end{equation}
%
where $\text{BR}_{it}^{\text{base}}$ is drawn from the MVN
(with no BM-BR correlation), $B_i^{*} = (B_i - \mu_B) /
\sigma_B$ is the standardized biomarker, $D_{it}$ is a binary
drug indicator (1 if on drug, 0 otherwise), and
$\sigma_{\text{BR}}$ is the BR standard deviation from the
response parameters.

The interaction is \emph{deterministic}: given a participant's
biomarker value and treatment status, the BR shift is fully
determined. The signal lives in the \emph{first moment}
(conditional mean).

\subsection{Why the distinction matters}

Under drug carryover, the drug effect persists into off-drug
periods and gradually decays. Under Architecture~B, the
correlation between biomarker and BR during off-drug periods
also decays (Equation~\ref{eq:mvn-cor}), eroding the signal
that the analysis model uses to detect the interaction. Under
Architecture~A, the moderation applies only during on-drug
timepoints and is absent during off-drug timepoints regardless
of carryover---the signal is binary (present or absent), not
graded. This architectural difference produces qualitatively
different conclusions about statistical power under carryover
(see the companion white paper,
\texttt{02-dgp-mean-moderation-vs-mvn.md}).

\section{The calibration problem}

Before describing the implementation attempts, we establish
the calibration target. For the two architectures to be
comparable, they should produce similar effect sizes for the
same value of $c_{bm}$. Under Architecture~B, the conditional
expectation (Equation~\ref{eq:cond-exp}) tells us that a
one-standard-deviation shift in the biomarker produces a BR
shift of:
%
\begin{equation}\label{eq:target}
  \Delta_{\text{BR}} = c_{bm} \cdot \sigma_{\text{BR}}
\end{equation}
%
For the parameters used in the Hendrickson et al.\ (2020)
framework ($c_{bm} = 0.45$, $\sigma_{\text{BR}} = 8$), this
is $0.45 \times 8 = 3.6$ units of BR per SD of biomarker.
Any correct Architecture~A implementation should produce a
shift of approximately the same magnitude.

We can verify this empirically. Drawing 500 participants
under Architecture~B with $c_{bm} = 0.45$, the regression
of mean on-drug BR on the biomarker yields a slope
corresponding to 3.48 BR units per SD of biomarker, with
residual SD of 3.89---consistent with the theoretical
prediction.

\section{Four implementation attempts}

\subsection{Attempt 1: multiplicative scaling}

The most intuitive approach is to scale each participant's
BR by a factor that depends on their biomarker value:
%
\begin{equation}
  \text{BR}_{it}^{\text{mod}} = \text{BR}_{it}^{\text{base}}
  \cdot (1 + \beta_{bm} \cdot B_i^{*})
\end{equation}

\begin{lstlisting}
dat[, (br_col) := get(br_col) * (1 + beta_bm * bm_centered)]
\end{lstlisting}

\textbf{Why it fails.} During off-drug timepoints, the base
BR value $\text{BR}_{it}^{\text{base}}$ is near zero (the
Gompertz response function evaluates to approximately zero
when time-on-drug is zero). Multiplying a near-zero value by
$(1 + \text{small})$ produces a near-zero result. The
moderation signal is proportional to the \emph{magnitude} of
the base BR, not to the drug exposure status. Consequently,
the analysis model sees strong biomarker-BR association during
on-drug periods (where base BR is large) and essentially no
association during off-drug periods---precisely mimicking
Architecture~B's correlation erosion.

\textbf{Result.} Power dropped approximately 60\% under
carryover, indistinguishable from Architecture~B.

\textbf{Lesson.} Multiplicative moderation conflates the
interaction signal with the response magnitude. When the
response is near zero (off-drug), the signal vanishes
regardless of the moderation strength.

\subsection{Attempt 2: additive with drug exposure decay}

To address the zero-base problem, we switched to an additive
formulation where the moderation signal tracks the
pharmacokinetic drug exposure level:
%
\begin{equation}
  \text{BR}_{it}^{\text{mod}} = \text{BR}_{it}^{\text{base}}
  + \beta_{bm} \cdot B_i^{*} \cdot \phi(t) \cdot \sigma_{\text{BR}}
\end{equation}
%
where the drug exposure function is:
%
\begin{equation}
  \phi(t) =
  \begin{cases}
    1 & \text{if on drug} \\
    \left(\tfrac{1}{2}\right)^{t_{sd} / t_{1/2}} &
      \text{if off drug with carryover} \\
    0 & \text{if never treated}
  \end{cases}
\end{equation}

\begin{lstlisting}
drug_exposure <- (1/2)^(tsd / t1half)
dat[, (br_col) := get(br_col) +
  beta_bm * bm_centered * drug_exposure * br_sd]
\end{lstlisting}

\textbf{Why it fails.} The drug exposure function decays
exponentially during off-drug periods. For the crossover
design with $t_{1/2} = 1.0$ week, the drug exposure at the
four off-drug timepoints (each 2.5 weeks apart) is 0.18,
0.031, 0.006, and 0.001. The additive moderation signal at
these timepoints is $0.45 \times 1 \times 0.18 \times 8 =
0.65$ BR units at the first off-drug timepoint, falling to
$0.004$ by the fourth. Against per-timepoint noise of
$\sigma_{\text{BR}} = 8$, this signal is undetectable.

\textbf{Result.} Power under carryover remained low,
comparable to Architecture~B.

\textbf{Lesson.} Tracking drug exposure decay in the
moderation signal reintroduces the very phenomenon that
Architecture~A is intended to avoid. If the moderation
decays with the drug, the architecture is functionally
equivalent to Architecture~B from the analysis model's
perspective.

\subsection{Attempt 3: additive with time-on-drug}

We then implemented the formulation found in the original
2025 tidyverse codebase
(\texttt{pmsimstats2025/analysis/scripts/pm\_functions.R},
initial commit \texttt{957c8bd}):
%
\begin{equation}
  \text{BR}_{it}^{\text{mod}} = \text{BR}_{it}^{\text{base}}
  + B_i \cdot \text{tod}_t \cdot \beta_{bm}
\end{equation}
%
where $\text{tod}_t$ is the cumulative time on drug at
timepoint $t$ (in weeks), and $B_i$ is the \emph{raw}
(uncentered, unstandardized) biomarker value.

\begin{lstlisting}
dat[, (br_col) := get(br_col) + bm * tod * beta_bm]
\end{lstlisting}

\textbf{Why it fails.} The biomarker in the Hendrickson
framework has mean $\mu_B \approx 124$ and SD
$\sigma_B \approx 15$ (systolic blood pressure). The
cumulative time-on-drug reaches 10--20 weeks. The moderation
shift for a typical participant at the last on-drug timepoint
is:
%
\begin{equation}
  124 \times 10 \times 0.45 = 558 \text{ BR units}
\end{equation}
%
The BR noise SD is 8. The signal-to-noise ratio is
approximately 70:1. Statistical power is 1.00 under all
conditions---not because Architecture~A is resilient to
carryover, but because the signal is so absurdly large that
no amount of noise or carryover could obscure it.

\textbf{Result.} Power = 1.00 everywhere. This was the
behavior observed in the original 2025 codebase, which was
incorrectly interpreted as carryover resilience.

\textbf{Lesson.} Dimensional analysis would have caught this
immediately. The product $B_i \times \text{tod}_t$ has units
of (mmHg $\times$ weeks), which is not commensurate with the
BR component (unitless symptom-score change). The resulting
signal is orders of magnitude larger than the noise floor.
Always check that the moderation term has the correct units
and magnitude relative to the outcome variable.

\subsection{Attempt 4: centered, binary, calibrated (correct)}

The correct formulation standardizes the biomarker, uses a
binary drug indicator, and scales by $\sigma_{\text{BR}}$ to
match the conditional expectation from the MVN model
(Equation~\ref{eq:cond-exp}):
%
\begin{equation}\label{eq:correct}
  \text{BR}_{it}^{\text{mod}} = \text{BR}_{it}^{\text{base}}
  + c_{bm} \cdot \frac{B_i - \mu_B}{\sigma_B}
    \cdot \mathbf{1}[\text{on drug at } t]
    \cdot \sigma_{\text{BR}}
\end{equation}

\begin{lstlisting}
bm_z <- (dat$bm - bm_mean) / bm_sd
for (tp in 1:nP) {
  if (d[tp]$onDrug) {
    br_col <- paste(tnames[tp], "br", sep = ".")
    dat[, (br_col) := get(br_col) + c.bm * bm_z * br_sd]
  }
}
\end{lstlisting}

\textbf{Why it works.}

\begin{enumerate}[leftmargin=2em]

\item \textbf{Centering.} The standardized biomarker
  $B_i^{*}$ has mean zero. The population-average moderation
  is zero---only deviations from the mean biomarker produce
  differential BR. This is the correct behavior: we are
  modeling an \emph{interaction} (differential response by
  biomarker level), not a main effect.

\item \textbf{Binary indicator.} The on-drug indicator
  $\mathbf{1}[\text{on drug}]$ is 1 during treatment and 0
  during off-drug periods. Unlike $\text{tod}$ (which grows
  with cumulative exposure) or $\phi(t)$ (which decays during
  washout), this indicator produces a clean, constant signal
  during all on-drug timepoints and zero signal during all
  off-drug timepoints. Carryover in the DGP affects the BR
  \emph{means} but not the moderation term.

\item \textbf{Scaling.} The factor $\sigma_{\text{BR}}$
  ensures that $c_{bm}$ has the same interpretation in both
  architectures: a one-SD shift in the biomarker produces a
  $c_{bm} \times \sigma_{\text{BR}}$ shift in BR. For
  $c_{bm} = 0.45$ and $\sigma_{\text{BR}} = 8$, this is
  3.6~BR units---matching the Architecture~B calibration
  target from Equation~\ref{eq:target}.

\end{enumerate}

\section{Empirical verification}

\subsection{Calibration}

Drawing 500 participants under both architectures with
identical parameters ($c_{bm} = 0.45$, $N = 500$, crossover
design, no carryover), the regression of mean on-drug BR on
the biomarker yields:

\begin{center}
\begin{tabular}{lcc}
\toprule
& \textbf{Arch A} & \textbf{Arch B} \\
\midrule
BR shift per SD of biomarker & 3.36 & 3.48 \\
Residual SD & 5.23 & 3.89 \\
\bottomrule
\end{tabular}
\end{center}

The shift magnitudes are comparable. Architecture~A has
slightly higher residual variance because the BM-BR
correlation in the covariance matrix is absent (the noise
components are independent), whereas Architecture~B's MVN
draw produces correlated noise that reduces the residual
variance in the regression.

\subsection{Power under carryover}

Using the N-of-1 (Hybrid) design with $N = 70$,
$c_{bm} = 0.45$, and 50 replicates per cell:

\begin{center}
\begin{tabular}{cccc}
\toprule
$t_{1/2}$ (weeks) & \textbf{Arch A} & \textbf{Arch B}
  & \textbf{Relative change (B)} \\
\midrule
0.0 & 0.74 & 0.82 & --- \\
0.5 & 0.68 & 0.64 & $-22\%$ \\
1.0 & 0.72 & 0.50 & $-39\%$ \\
\bottomrule
\end{tabular}
\end{center}

Architecture~A power is stable (variation within Monte Carlo
error at $N_{\text{reps}} = 50$). Architecture~B shows
progressive power loss with increasing carryover half-life,
reaching a 39\% relative decline at $t_{1/2} = 1.0$ week.

\section{Discussion}

\subsection{Why scaling matters}

The four attempts illustrate a general principle in
simulation development: the \emph{units and magnitude} of
each component in the data-generating process must be
verified against a known reference. Three of the four
attempts produced code that executed without errors,
generated plausible-looking data, and could be analyzed by
the downstream statistical model. The errors were detectable
only by checking whether the resulting power estimates were
consistent with theoretical expectations.

\begin{itemize}[leftmargin=2em]

\item Attempt~1 was identifiable by checking the correlation
  between biomarker and BR during off-drug timepoints
  (expected: nonzero under Architecture~A; observed: zero).

\item Attempt~2 was identifiable by computing the
  signal-to-noise ratio at off-drug timepoints (expected:
  detectable; observed: negligible).

\item Attempt~3 was identifiable by dimensional analysis
  (biomarker in mmHg, time in weeks, product in
  mmHg$\cdot$weeks, outcome in symptom-score units) or by
  checking the signal-to-noise ratio (expected: $\sim$1;
  observed: $\sim$70).

\end{itemize}

\subsection{The role of the conditional expectation}

The correct formulation (Equation~\ref{eq:correct}) was
derived by asking: ``What effect size does Architecture~B
produce for the same value of $c_{bm}$?'' The answer comes
directly from the conditional expectation of the MVN
distribution (Equation~\ref{eq:cond-exp}). By matching this
conditional expectation, we ensure that $c_{bm} = 0.45$
means the same thing---in terms of detectable effect
size---under both architectures. This is essential for the
comparative power analysis that motivates the dual-architecture
framework.

\subsection{Implementation}

The \texttt{dgp\_architecture} parameter controls the choice
between architectures in the \texttt{pmsimstats-ng} R
package. It is threaded through three functions:

\begin{itemize}[leftmargin=2em]

\item \texttt{buildSigma()}: When
  \texttt{dgp\_architecture = "mean\_moderation"}, the BM-BR
  correlation block is omitted from the covariance matrix.
  The biomarker and BR components are drawn independently.

\item \texttt{generateData()}: After the MVN draw, the
  additive moderation term
  (Equation~\ref{eq:correct}) is applied to all on-drug BR
  columns.

\item \texttt{generateSimulatedResults()}: The parameter is
  passed to the sigma cache and all \texttt{generateData()}
  calls within the simulation loop.

\end{itemize}

The default is \texttt{"mvn"} (Architecture~B), preserving
backward compatibility with the Hendrickson et al.\ (2020)
publication code. All 67 existing unit tests pass under the
default.

\vfill

\noindent\rule{\textwidth}{0.4pt}

{\footnotesize
Rendered on 2026-03-25 at 18:49 PDT.\\
Source: \textasciitilde/prj/alz/10-pmsimstats-ng/pmsimstats-ng/docs/19-mean-moderation-implementation-notes.tex}

\end{document}
